% !PRINTED 55
\weberpage{55}

%<<CONT
piece of surface, where zero points and points of discontinuity of order $m$ are to be counted as $m$ such points of the first order.
From this it follows that in the interior of a finite piece of surface there also lies only a finite number of zero points.

In the same sense there results

7. The integral
\[
\frac{1}{2\pi i}\int u\,d\log\varphi(u)
\]
is equal to the sum of the values of $u$ for which $\varphi(u)$ vanishes, diminished by the sum of the values of $u$ for which $\varphi(u)$ becomes infinite.

8. A function which possesses no points of discontinuity at all in the finite region is called an \gesperrt{integral function}. \gesperrt{An integral function which also remains finite at infinity is a constant.} An integral function which also has no essentially singular point at infinity, for which there is therefore an exponent $m$ such that $u^{-m}\varphi(u)$ does not become infinite for $u = \infty$, is a rational function.

9. An integral function whose reciprocal is likewise integral is called a \gesperrt{unit function}. Such a function becomes neither infinite nor zero in the finite region. Its logarithm is therefore likewise an integral function, and from this it follows that every unit function can be represented in the form $e^{g(u)}$, in which $g(u)$ is an integral function. Conversely, every expression of this form is a unit function.

10. According to the fundamental investigations of \gesperrt{Weierstrass} on single-valued analytic functions (Abhandlungen der Berliner Akademie von 1876) there is always an integral function $G(u)$ which vanishes at the points of discontinuity of a function $\varphi(u)$ and only at these, and this function $G(u)$ is determined by the points of discontinuity of $\varphi(u)$ up to a unit function as a factor. Then $\varphi(u)G(u) = G_1(u)$ is likewise an integral function, and one can therefore represent every function $\varphi(u)$ as the quotient of two integral functions
\[
\varphi(u) = \frac{G_1(u)}{G(u)}
\]
in which $G(u)$ and $G_1(u)$ have no zero points in common.
